\chapter{序論}

序論では，本研究の背景である近年のSNSの動向，Twitterの概要及び問題について述べた上で，本研究の目的と構成について説明する．

\section{背景}
\subsection{近年のSNSの動向}
株式会社ICT総研の「2020年度 SNS利用動向に関する調査」 \cite{ICT総研}によると日本国内におけるSNS (Social Networking Service)の利用者は年々増加の傾向にあり，2020年末には7975万人に達する見込みである(図\ref{日本におけるSNS利用者数})．これはネット利用人口全体の80.3\%がSNSを利用している計算となり，SNSが社会における重要なインフラであることを示している．さらに，このまま普及が進めば2022年末には利用者数は8241万人，ネットユーザー全体に占める利用率は83.3\%に達する見通しである．

\begin{figure}[htbp]
	\begin{center}
		\includegraphics[width=11cm]{./figure/fig1.png}
		\caption{日本におけるSNS利用者数}
		\label{日本におけるSNS利用者数}
	\end{center}
\end{figure}
\newpage
また，株式会社ICT総研が2020年7月に実施したアンケート調査\cite{ICT総研}では，4400人のアンケート対象者のうち8割以上の人が「人とのコミュニケーション」を行う目的でSNSや通話・メールアプリを利用していると回答している．そして図\ref{主なSNSの利用率}のように全サービスの中で最も利用者が多かったのがLINEで77.4\%，そのあとに続いたのがTwitterの38.5\%，Instagramの35.7\%，YouTubeの23.2\%，Facebookの21.7\%，TikTokの8.1\%，Skypeの7.1\%である．本研究ではこれらのSNSのうちTwitterに焦点を当てる．
\begin{figure}[htbp]
	\begin{center}
		\includegraphics[width=14cm]{./figure/fig2.png}
		\caption{主なSNSの利用率}
		\label{主なSNSの利用率}
	\end{center}
\end{figure}

\subsection{Twitterの概要及び問題}
Twitter \cite{Twitter}は米国Twitter社が運営するSNSで，ユーザは「ツイート」と呼ばれる半角280文字（日本語，中国語，韓国語は全角140文字）以内のテキストや画像，動画，URLを投稿することができる．Twitterではあるユーザから他のユーザのアカウントをフォローすることで，そのアカウントのツイートを見ることができる．Twitterの用途は人それぞれであり，例えば，あるユーザは親しいユーザをフォローし，そのユーザとツイートやダイレクトメッセージなどの機能を使用してコミュニケーションを行ったり，またあるユーザは自分の興味がある分野に関する情報を積極的にツイートするアカウントをフォローし，情報収集を行っている．さらに，Twitterには速報性があるのでユーザはタイムリーに社会的な関心事や自分の興味に基づいた情報を入手できるという特徴もある．しかし，自分が興味ある内容だからといって，必ずしも今必要な情報だとは限らない．その代表例が「ネタバレ」である．

\subsection{ネタバレ}
%ネタバレとは映画，劇，小説，劇，漫画，アニメなどの作品やスポーツの試合についての内容を知っている者が，知らない者に対してその内容を暴露してしまうこと，またその情報のことである．
中村ら \cite{中村}によれば「ネタバレ」とは，映画やテレビドラマ，スポーツ等，ある対象コンテンツをまだ楽しんだことがなく，これから楽しもうと思っているユーザに，その対象に関する経過や結果を知らせること，またその情報のことである．例えば，スポーツ中継をリアルタイムで視聴できないため，録画をして後で視聴する人もいる．%そのような録画視聴を楽しみにしている視聴者が，視聴前にそのスポーツの結果に関連するツイートを見てしまうとそのツイートは「ネタバレ」となり，コンテンツの面白さを低下させる恐れがある．
そのような視聴者が，視聴前にそのスポーツの結果に関連するネタバレツイートを見てしまうと，試合の面白さを低下させる恐れがある．
実際，白鳥ら \cite{白鳥1}はネタバレが視聴者の試合に対する緊張感や一喜一憂度合いを減少させることを証明している．本研究ではネタバレの中でも欧州4大リーグのサッカーの試合のネタバレに着目する．サッカーを研究対象とした理由は，サッカーが試合進行に関するイベントが任意のタイミングで生じ得るスポーツだからである．すなわち，野球，テニス，バレーボール等の試合進行がターン制のスポーツと異なり，サッカーの試合はどのタイミングで試合進行に関するツイートが生じ得るかまったく予想が出来ず，ネタバレが視聴者に与える影響が大きいと考えられるからである．

\subsection{欧州4大リーグ}
欧州4大リーグとは，イングランド1部のプレミアリーグ，イタリア1部のセリエA，スペイン1部のラ・リーガ，ドイツ1部のブンデスリーガの４つのリーグを総称したもので，これら4つにフランス1部のリーグ・アンを加えて欧州５大リーグとも呼ばれる．欧州サッカー連盟(UEFA)に所属する国のサッカーリーグについて公式戦での成績をもとに算出される「UEFA Country coefficients」 \cite{UEFA}というランキングでは，1位がラ・リーガ，2位がプレミアリーグ，3位がセリエA，4位がブンデスリーガとなっており，4大リーグの競技レベルの高さが窺える．

\section{本研究の目的}
本章では，日本におけるSNSの現状，Twitterにおけるネタバレ問題について述べた．そこで本研究では，Twitter上のネタバレ，特にサッカーの試合に関するネタバレを自動で検出し閲覧を防止する機械学習モデルを構築することを目的とする．さらに，ネタバレは試合の重要な情報を含んでいると考えられるため，検出するネタバレをスコアに関係するネタバレとスコアに関係しないネタバレに分類する．これによりツイートから試合の重要な情報を抽出することが可能になり，本モデルは，ネタバレ防止の他に，試合速報やハイライトなど試合情報の自動要約に応用されることが見込まれる．

\section{本論文の構成}
本章では，本研究の背景及び目的を述べた．第2章では，ネタバレの影響や検出に関連する研究を紹介する．第3章では提案手法，第4章では具体的な実装方法について述べる．第5章で実験の結果を示し，第6章で結果の考察，第7章で本論文のまとめ及び今後の展望について述べる．

%本論文の流れは以下に示す．
%\begin{enumerate}
%   \item 第1章　序論
%	　　本研究の背景と目的，そして本論文の構成について述べる．近年のSNSの現状について述べ，問題提起を行う．そして本研究が提案する手法の
 %  \item 第2章　関連研究
%	　　あ
%   \item 第3章　提案手法
%	　　あ
%   \item 第4章　実装
%	　　あ
%   \item 第5章　結果
%	　　あ
%   \item 第6章　考察
%	　　あ
 %  \item 第7章　結論
%	　　あ
%\end{enumerate}
